\clearpage
\section{Additional Figures and Tables}\label{app:additional}

This appendix collects supporting results from the project's notebooks 02--07, in the order of
the analysis: the event study (notebook 02), state dependence (03), the walk-forward CPI strategy
(04), the broad surprise book (05), surprise and DTW (06) and the multi-market book (07). Unless noted, the evaluation window is 2019--2026, the cost is two
ticks plus \$4.50 per round trip, and every estimate uses only earlier data.

\FloatBarrier
\subsection{Event study (notebook 02)}

Figure~\ref{fig:app_es_event} shows, for \ES, how large the move after CPI, payrolls and PCE is at
each horizon and how much of it the surprise explains. The move after CPI keeps growing for about
30 minutes; the share explained by the surprise is highest at one minute and stays near 0.25,
while payrolls move \ES\ almost as much but with an $R^2$ below 0.07. Figure~\ref{fig:app_es_beta}
shows the coefficient of each headline surprise by horizon: a one-standard-deviation hot CPI
print lowers \ES\ by 18 bps after one minute and by about 24 bps after 30 minutes. Table~\ref{tab:app_impact}
reports the mean absolute move in all four markets.

\begin{figure}[!htb]
  \centering
  \begin{subfigure}[b]{0.49\linewidth}\centering
    \includegraphics[width=\linewidth]{app_macro_market_impact}
    \caption{Mean absolute \ES\ return after each release.}
  \end{subfigure}\hfill
  \begin{subfigure}[b]{0.49\linewidth}\centering
    \includegraphics[width=\linewidth]{app_macro_r2_comparison}
    \caption{$R^2$ of the surprise regressions.}
  \end{subfigure}
  \caption{\ES\ after CPI, payrolls (NFP) and PCE, 2016--2026, by horizon.}
  \label{fig:app_es_event}
\end{figure}

\begin{figure}[!htb]
  \centering
  \includegraphics[width=0.62\linewidth]{app_macro_surprise_coefficients}
  \caption{\ES\ response to a one-standard-deviation surprise in each headline statistic, by
  horizon (bps), from regressions on all surprise components of the release.}
  \label{fig:app_es_beta}
\end{figure}

\begin{table}[!htb]
  \centering\small
  \caption{Mean absolute return after the release (bps), 2016--2026.}\label{tab:app_impact}
  \begin{tabular}{llccccc}
    \toprule
    Market & Release & 1 min & 5 min & 15 min & 30 min & 60 min \\
    \midrule
    \ES & CPI & 28.1 & 31.4 & 34.6 & 38.5 & 37.8 \\
        & NFP & 21.5 & 26.3 & 26.1 & 29.6 & 29.2 \\
        & PCE &  7.0 &  9.4 & 12.9 & 14.9 & 18.1 \\
    \NQ & CPI & 38.6 & 44.7 & 48.4 & 53.6 & 54.0 \\
        & NFP & 25.6 & 31.6 & 31.5 & 36.0 & 35.1 \\
        & PCE &  9.2 & 12.5 & 17.2 & 18.5 & 23.2 \\
    \YM & CPI & 23.7 & 27.7 & 29.6 & 33.2 & 33.3 \\
        & NFP & 19.3 & 23.7 & 25.0 & 27.9 & 28.6 \\
        & PCE &  5.4 &  7.3 & 10.2 & 11.9 & 15.9 \\
    \ZN & CPI & 18.4 & 20.2 & 21.4 & 22.5 & 24.0 \\
        & NFP & 20.0 & 21.6 & 22.2 & 23.6 & 24.4 \\
        & PCE &  4.9 &  6.1 &  7.3 &  8.3 &  9.3 \\
    \bottomrule
  \end{tabular}
\end{table}

\FloatBarrier
\subsection{State dependence (notebook 03)}

Table~\ref{tab:app_states} lists the six pre-registered state variables, all measured before the
release, with the hypothesis each was chosen to test. Four further variables are reported in the
figures but were not used to select models: 5-day momentum, the gap to the 50-day moving average,
realized volatility over the last 60 minutes, and the return over the last 60 minutes.

\begin{table}[!htb]
  \centering\small
  \caption{Pre-registered state variables and hypotheses.}\label{tab:app_states}
  \begin{tabular}{p{0.15\linewidth}p{0.36\linewidth}p{0.41\linewidth}}
    \toprule
    State & Definition (known before $T$) & Hypothesis \\
    \midrule
    \texttt{rv\_20d} & 20-day realized volatility of daily closes (annualized \%) & High-volatility regimes give a larger surprise beta (uncertainty, thinner liquidity) \\
    \texttt{rv\_pre\_1d} & Realized volatility of 5-minute returns over the last $\sim$23 trading hours (bps) & Short-term stress amplifies the reaction \\
    \texttt{mom\_20d} & 20-day log return (\%) & After strong rallies, bad news hits harder (stretched positioning) \\
    \texttt{dd\_60d} & Distance from the 60-day high (\%) & In drawdowns the market is more sensitive to inflation news \\
    \texttt{overnight\_ret} & Return from the prior 16:00 ET close to $T-1$ (bps) & A large pre-release move signals partial pre-positioning, so less is left to react \\
    \texttt{prev\_x} & Previous signed surprise of the same family & Surprise streaks: a second hot print confirms the trend, so the response is larger \\
    \bottomrule
  \end{tabular}
  \par\smallskip\footnotesize Exploratory (reported, not used for selection): \texttt{mom\_5d},
  \texttt{ma50\_gap}, \texttt{rv\_pre\_60m}, \texttt{ret\_pre\_60m}. Each state enters
  equation~\eqref{eq:state} as a percentile against its trailing one-year distribution.
\end{table}

Figure~\ref{fig:app_daily_state} shows three of the daily states in each market from 2010 to
2026, computed on roll-adjusted daily closes. The states vary a lot over the sample, so the terciles
used below compare genuinely different regimes: 20-day volatility in the equity indices is near
10\% in the quiet years before 2018 and peaks at 90--100\% in March 2020, with further spikes in
2011, 2015, 2018, 2022 and April 2025, and drawdowns from the 60-day high reach $-35$ to $-40\%$ in
2020. \ZN's states move with the rate cycle rather than with equity stress: its volatility is
elevated through 2022--2023 and its deepest drawdown ($-10\%$) is in 2022. The \ES\ volatility of
about 85\% in 2011 is overstated: the notebook computes it from the original vendor file, whose
early sessions are incomplete; \YM, from a complete file, peaks at about 45\% in the same episode.

\begin{figure}[!htb]
  \centering
  \includegraphics[width=\linewidth]{app_daily_state_4markets}
  \caption{Daily state features, 2010--2026: 20-day realized volatility (annualized \%), 20-day
  momentum (\%) and distance from the 60-day high (\%), computed on roll-adjusted daily closes. The
  2011 \ES\ volatility spike reflects incomplete sessions in the original vendor file.}
  \label{fig:app_daily_state}
\end{figure}

Figure~\ref{fig:app_buckets} splits the five-minute CPI beta by tercile of each pre-release state
in \ES\ and \YM. The beta rises from under 10 bps per standard deviation in the calmest third of
20-day volatility to 40--70 bps in the middle and top thirds in both markets, the clearest pattern among the
six states. Table~\ref{tab:app_state} lists the strongest of the 48 primary tests in
each market. Figure~\ref{fig:app_pvalues} shows that the analytic HC3 and permutation $p$-values
agree closely, so the inference does not depend on the normal approximation.

\begin{figure}[!htb]
  \centering
  \includegraphics[width=\linewidth]{app_bucket_betas_ES_YM}
  \caption{Five-minute CPI surprise beta (bps per SD, 95\% intervals) by tercile of each
  pre-release state: 20-day realized volatility, previous-day volatility, 20-day momentum, 60-day
  drawdown, overnight return and previous surprise.}
  \label{fig:app_buckets}
\end{figure}

\begin{table}[!htb]
  \centering\small
  \caption{Strongest state interactions among the 48 primary tests, 2016--2026.}\label{tab:app_state}
  \begin{threeparttable}
  \begin{tabular}{lllcccc}
    \toprule
    Market & Interaction & Target & $d$ (bps) & $t$ (HC3) & $p$ (perm.) & $q$ (BH) \\
    \midrule
    \ES & CPI $\times$ 20-day vol & 5-min reaction & $+61.1$ & 3.81 & 0.000 & 0.010 \\
        & NFP $\times$ 20-day vol & 30-min drift & $-12.2$ & $-2.99$ & 0.016 & 0.082 \\
        & All $\times$ overnight return & 30-min drift & $+7.4$ & 2.65 & 0.021 & 0.134 \\
    \YM & CPI $\times$ 20-day vol & 5-min reaction & $+58.7$ & 3.74 & 0.004 & 0.014 \\
        & NFP $\times$ previous-day vol & 5-min reaction & $-21.0$ & $-3.07$ & 0.012 & 0.042 \\
        & NFP $\times$ 20-day vol & 5-min reaction & $-17.3$ & $-3.11$ & 0.014 & 0.042 \\
    \bottomrule
  \end{tabular}
  \begin{tablenotes}\footnotesize
    \item Notes: $d$ is the change in the surprise beta from the calmest to the most stressed state
    in equation~\eqref{eq:state}. $q$: Benjamini--Hochberg over the 48 primary tests per market. The
    NFP effects in \YM\ come entirely from 2016--2020.
  \end{tablenotes}
  \end{threeparttable}
\end{table}

\begin{figure}[!htb]
  \centering
  \includegraphics[width=\linewidth]{app_fig13_pvalue_comparison_3markets}
  \caption{Analytic (HC3) against permutation $p$-values for the primary state tests, \ES, \NQ\ and \ZN.}
  \label{fig:app_pvalues}
\end{figure}

\paragraph{Interaction regressions (notebook 03, section 9).} Rather than splitting the releases
into terciles, the interaction regression uses the full sample at once. For release $i$,
\begin{equation}\label{eq:app_inter}
  r_i = a + b\,x_i + c\,s_i + d\,(x_i\times s_i) + \varepsilon_i,\qquad s_i=\text{pct}_i-0.5,
\end{equation}
where $x_i$ is the signed surprise and $\text{pct}_i$ the percentile of the state against its
trailing one-year distribution. Centering the state makes $b$ the surprise beta at the median
state; $c$ asks whether the state predicts the return on its own; and $d$, the key coefficient, is
the change in the surprise beta from the lowest to the highest state percentile. Because $x$ is
signed so that $b>0$, a positive $d$ means a stronger reaction in the high state. The pooled model
(\texttt{ALL}) gives each family its own intercept and slope but a common $c$ and $d$. Each market
has 360 regressions: ten states (six primary, four exploratory), four families (CPI, NFP, PCE and
pooled) and nine targets (the reaction from the pre-release price at 1--60 minutes and the drift
from one minute after the release at 5--60 minutes). Standard errors are HC3. The 48 primary tests
(six states, four families, two targets: the five-minute reaction and the 30-minute drift) also get
a 2,000-draw permutation test and the Benjamini--Hochberg correction. A credible state effect looks
like a row of same-colored cells across horizons; a single bold cell among about 90 is what chance
produces, since at 5\% significance four to five false positives are expected per heat map.

The CPI heat maps are in Figure~\ref{fig:state} of the main text. Figures~\ref{fig:app_inter_nfp}--\ref{fig:app_inter_all}
show the other families. For payrolls (Figure~\ref{fig:app_inter_nfp}), the volatility rows are
negative: strong payrolls help equities less when volatility is high. The pattern is strongest in
\YM, where both 20-day and previous-day volatility have $t$-statistics of $-3.1$ to $-4.5$ on the
5- to 60-minute reactions, and present in \ES\ on the drifts; but in \YM\ it comes entirely from 2016--2020,
so it is not used for trading. PCE (Figure~\ref{fig:app_inter_pce}) shows only scattered cells with
no consistent row, as expected for a release whose surprise explains little of the move. In the
pooled model (Figure~\ref{fig:app_inter_all}), the overnight return and 5-day momentum rows are
positive on the 15--60-minute drifts in \ES\ and \YM\ ($t$ of 2.1--2.9): the drift follows the surprise
more strongly after a market has already been rising. Only the \ES\ overnight-return effect on the
30-minute drift is a primary test, and it does not survive the correction ($q=0.13$). In \ZN, the
previous-day volatility row is negative on every drift ($t$ of $-2.1$ to $-3.2$): after volatile
days, the little drift \ZN\ has is even smaller.

\begin{figure}[!htb]
  \centering
  \includegraphics[width=\linewidth]{app_interaction_NFP_4markets}
  \caption{Payrolls (NFP): $t$-statistics of the surprise $\times$ state interaction $d$ in
  equation~\eqref{eq:app_inter}, by state (rows) and target (columns). Bold: $|t|>1.96$.}
  \label{fig:app_inter_nfp}
\end{figure}

\begin{figure}[!htb]
  \centering
  \includegraphics[width=\linewidth]{app_interaction_PCE_4markets}
  \caption{PCE: $t$-statistics of the surprise $\times$ state interaction $d$, by state and target.}
  \label{fig:app_inter_pce}
\end{figure}

\begin{figure}[!htb]
  \centering
  \includegraphics[width=\linewidth]{app_interaction_ALL_4markets}
  \caption{All three families pooled (family-specific intercepts and slopes, common $c$ and $d$):
  $t$-statistics of the surprise $\times$ state interaction, by state and target.}
  \label{fig:app_inter_all}
\end{figure}

\FloatBarrier
\subsection{Walk-forward CPI strategy (notebook 04)}

Figure~\ref{fig:app_latency} shows how the CPI strategy's Sharpe ratio changes with the entry
delay and the exit time. In both \ES\ and \YM\ the 30-minute exit is the best column at every
delay, and short holds lose most of their value when the entry is delayed. Figure~\ref{fig:app_paths}
shows the walk-forward estimates behind the state model: the CPI drift slope in \ES\
fell to about 1 bp per standard deviation in 2021 and rose to 6--8 bps from 2023, and the
$t$-statistic of the volatility interaction stays below about 2 until 2022 in \ES\ and \NQ\ and
never reaches it in \ZN.
Figures~\ref{fig:app_peer}--\ref{fig:app_peer_all} show the walk-forward strategies in the peer
chart format (notebook 04, section 16.5): the strategy before and after costs against being
always long, one contract, in exactly the same 30-minute windows. In the three equity indices the
strategies and the always-long line diverge from 2022: trading the CPI surprise gains 6.6\% in
\ES, 7.5\% in \NQ\ and 5.0\% in \YM\ while always long in the same windows loses 5--6\%, so the
profit comes from the direction of the surprise, not from being in the market at release times.
The naive sign rule and the model on CPI, PCE and payrolls show the same pattern. In \ZN\ the
surprise model trades only a handful of releases, because its predicted drift rarely exceeds the
cost of one \ZN\ tick, and the naive rule loses after costs.

\begin{figure}[!htb]
  \centering
  \includegraphics[width=0.85\linewidth]{app_latency_grid_ES_YM}
  \caption{Net Sharpe ratio of the CPI strategies by extra entry delay (rows) and exit time
  (columns), 2019--2026. Left: surprise only; right: surprise and state.}
  \label{fig:app_latency}
\end{figure}

\begin{figure}[!htb]
  \centering
  \includegraphics[width=0.62\linewidth]{app_fig17_coefficient_paths_3markets}
  \caption{Walk-forward estimates for CPI, re-estimated before each release: the surprise slope and
  the surprise $\times$ 20-day volatility interaction with their $t$-statistics, \ES, \NQ\ and \ZN.}
  \label{fig:app_paths}
\end{figure}

\begin{figure}[!htb]
  \centering
  \includegraphics[width=\linewidth]{app_peer_CPI_surprise_only_4markets}
  \caption{Peer chart format, CPI surprise model, 2019--2026: cumulative return of one contract,
  gross and net of costs, against always long in the same windows.}
  \label{fig:app_peer}
\end{figure}

\begin{figure}[!htb]
  \centering
  \includegraphics[width=\linewidth]{app_peer_CPI_naive_4markets}
  \caption{Peer chart format, CPI naive sign rule (every CPI release, sign of the headline
  surprise), 2019--2026.}
  \label{fig:app_peer_naive}
\end{figure}

\begin{figure}[!htb]
  \centering
  \includegraphics[width=\linewidth]{app_peer_ALL_surprise_only_4markets}
  \caption{Peer chart format, surprise model on CPI, PCE and payrolls, 2019--2026.}
  \label{fig:app_peer_all}
\end{figure}

\paragraph{Final comparison.} Table~\ref{tab:app_final} and Figure~\ref{fig:app_final} compare
the five CPI strategies of notebook 04 (section 14) in all four markets and three release
universes. In the equity indices every CPI strategy has a positive Sharpe ratio, and the naive
sign rule is the best or close to it (1.03 in \ES, 1.00 in \NQ, 0.99 in \YM); conditioning on the
pre-release state never improves on the surprise alone by a significant margin (the smallest
bootstrap probability that a state model is not better is 0.23). Adding PCE and payrolls lowers the
naive rule in \ES\ and \NQ\ but less so in \YM. In \ZN\ the fitted models almost never trade,
because their predicted drift rarely exceeds one tick, and every strategy is negative.

\begin{table}[!htb]
  \centering\small
  \caption{Walk-forward CPI strategies, net Sharpe ratio (90\% block-bootstrap interval), 2019--2026,
  2 ticks + \$4.50.}\label{tab:app_final}
  \begin{threeparttable}
  \setlength{\tabcolsep}{4pt}
  \begin{tabular}{lcccc}
    \toprule
    Strategy & \ES & \NQ & \YM & \ZN \\
    \midrule
    \multicolumn{5}{l}{\emph{CPI only}} \\
    Naive sign & 1.03 [0.32, 1.76] & 1.00 [0.37, 1.66] & 0.99 [0.27, 1.72] & $-0.05$ [$-0.82$, 0.56] \\
    Surprise only & 0.99 [0.25, 1.72] & 0.79 [0.12, 1.47] & 0.84 [0.07, 1.61] & $-0.51$ (2 trades) \\
    Surprise + State & 0.77 [0.01, 1.51] & 0.75 [0.04, 1.44] & 0.83 [0.13, 1.51] & $-0.40$ (5 trades) \\
    Surprise + State filter & 1.07 [0.40, 1.67] & 0.82 [0.20, 1.38] & 0.98 [0.25, 1.65] & $-0.37$ (1 trade) \\
    Surprise + Adaptive state & 0.95 [0.26, 1.61] & 0.75 [0.04, 1.44] & 0.74 [0.02, 1.44] & $-0.51$ (2 trades) \\
    \midrule
    \multicolumn{5}{l}{\emph{CPI, PCE and payrolls}} \\
    Naive sign & 0.33 & 0.38 & 0.57 & $-0.65$ \\
    Surprise only & 0.59 & 0.37 & 0.78 & $-0.22$ \\
    Surprise + State filter & 0.65 & 0.37 & 0.81 & $-0.07$ \\
    \midrule
    CPI trades, naive / surprise only & 86 / 52 & 79 / 62 & 77 / 48 & 86 / 2 \\
    \bottomrule
  \end{tabular}
  \begin{tablenotes}\footnotesize
    \item Notes: Entry at $T+1$, exit at $T+30$; every slope, sign and state model re-estimated on
    earlier releases only. Intervals from a three-month block bootstrap.
  \end{tablenotes}
  \end{threeparttable}
\end{table}

\begin{figure}[!htb]
  \centering
  \includegraphics[width=\linewidth]{app_final_sharpe_4markets}
  \caption{Final comparison of notebook 04: walk-forward net Sharpe ratio with 90\% bootstrap
  intervals by strategy and release universe (CPI; CPI + PCE; CPI + PCE + NFP), 2019--2026.}
  \label{fig:app_final}
\end{figure}

\FloatBarrier
\subsection{Broad surprise book (notebook 05)}

\paragraph{Which releases earn money? (notebook 05, section 9.2).} Table~\ref{tab:app_family} and
Figure~\ref{fig:app_family} split the broad book's net P\&L (naive sizing, 2019--2026) by release
family. CPI is the largest contributor in every equity index: $+4.3\%$ of notional in \ES, $+6.0\%$
in \NQ\ and $+3.8\%$ in \YM. Behind it the ranking differs by market. In \ES\ the next contributors are
JOLTS ($+3.0\%$), retail sales ($+2.3\%$) and the University of Michigan survey ($+1.4\%$); in \NQ\
factory orders ($+4.6\%$), the University of Michigan survey ($+4.2\%$), retail sales ($+2.6\%$) and
GDP ($+2.1\%$); in \YM\ ISM Manufacturing ($+3.3\%$), JOLTS ($+2.3\%$) and retail sales ($+1.7\%$).
Payrolls, the release with the largest market impact, earns nothing or loses in all three ($-1.7\%$,
$-0.5\%$ and $-0.1\%$), and ISM Services, the regional Fed surveys and pending home sales lose small
amounts. No major family is individually significant: the $t$-statistics of the main families are at
most 1.9 with naive sizing, and CPI with volatility sizing reaches 2.0 in \ES\ and 1.9 in \NQ. The broad book
therefore works through diversification across many small edges rather than through one release.
In \ZN\ every family with more than a handful of trades loses, led by payrolls ($-4.9\%$), retail
sales ($-4.5\%$) and CPI ($-3.0\%$): its hit rates of 22--47\% reflect a drift smaller than the cost.

\begin{table}[!htb]
  \centering\small
  \caption{Net P\&L of the broad surprise book by release family, 2019--2026: total net return
  (\% of notional) and $t$-statistic of the mean trade, naive sizing.}\label{tab:app_family}
  \begin{threeparttable}
  \setlength{\tabcolsep}{3pt}
  \begin{tabular}{lcclcclcc}
    \toprule
    \multicolumn{3}{c}{\ES} & \multicolumn{3}{c}{\NQ} & \multicolumn{3}{c}{\YM} \\
    \cmidrule(lr){1-3}\cmidrule(lr){4-6}\cmidrule(lr){7-9}
    Family & Total & $t$ & Family & Total & $t$ & Family & Total & $t$ \\
    \midrule
    CPI & 4.3\% & 1.40 & CPI & 6.0\% & 1.50 & CPI & 3.8\% & 1.48 \\
    JOLTS & 3.0\% & 1.56 & Factory Orders & 4.6\% & 1.89 & ISM Manuf. & 3.3\% & 1.02 \\
    Retail Sales & 2.3\% & 1.12 & U.\ of Michigan & 4.2\% & 1.45 & JOLTS & 2.3\% & 1.19 \\
    U.\ of Michigan & 1.4\% & 0.67 & Retail Sales & 2.6\% & 1.13 & Retail Sales & 1.7\% & 1.01 \\
    Adv.\ Goods Trade & 1.3\% & 1.13 & GDP & 2.1\% & 1.55 & Trade Balance & 1.2\% & 0.88 \\
    \midrule
    ISM Services & $-1.5\%$ & $-0.49$ & ISM Services & $-1.3\%$ & $-0.32$ & Wholesale Trade & $-1.6\%$ & $-2.40$ \\
    Payrolls (NFP) & $-1.7\%$ & $-0.60$ & Richmond Fed & $-1.9\%$ & $-1.02$ & Wholesale Inv. & $-2.7\%$ & $-1.68$ \\
    \bottomrule
  \end{tabular}
  \begin{tablenotes}\footnotesize
    \item Notes: Top five contributors and the two largest losers in each market. Payrolls: $-0.5\%$
    in \NQ\ and $-0.1\%$ in \YM. In \ZN\ every family with more than a few trades loses.
  \end{tablenotes}
  \end{threeparttable}
\end{table}

\begin{figure}[!htb]
  \centering
  \includegraphics[width=\linewidth]{app_fig26_family_contribution_3markets}\\[4pt]
  \includegraphics[width=0.42\linewidth]{app_nb05b_YM_family_contribution}
  \caption{Total net P\&L of the broad surprise book by release family, 2019--2026: \ES, \NQ\ and
  \ZN\ (top) and \YM\ (bottom).}
  \label{fig:app_family}
\end{figure}

\paragraph{CPI only against CPI plus other surprises (notebook 05, section 10).}
Figure~\ref{fig:app_nb05_final} and Table~\ref{tab:app_nb05_final} compare the CPI-only books (blue)
with the books that add the other screened releases (orange and green), against the DTW book from
the team's original slide (gray; \ES, 2018--2026, used as a fixed reference in every panel). The
trade-off is the same in every equity index. The CPI books earn the most per trade (5--16 bps net)
but trade only about 11 times a year. Adding the other releases raises the count to about 95 a year
at about 2 bps per trade, and the simple naive sign rule is then the best broad book (0.65 in \ES,
0.71 in \NQ, 0.68 in \YM), close to the CPI books' Sharpe ratio with far more trades, and it keeps the
highest Sharpe ratio without 2022 (0.28--0.44). The fitted surprise model and volatility scaling
help on CPI but not on the broad book, where scaling mainly adds drawdown ($-12$ to $-14\%$).
Trading a second 30-minute window doubles the trades and removes the edge. In \ZN\ every book
loses, and the fitted models almost never trade.

\begin{table}[!htb]
  \centering\small
  \caption{CPI only against CPI plus other surprises: net Sharpe ratio (trades per year),
  2019--2026, 2 ticks + \$4.50.}\label{tab:app_nb05_final}
  \begin{threeparttable}
  \setlength{\tabcolsep}{4pt}
  \begin{tabular}{lcccc}
    \toprule
    Book & \ES & \NQ & \YM & \ZN \\
    \midrule
    CPI $\cdot$ naive & 0.51 (12) & 0.54 (12) & 0.54 (11) & $-0.67$ (12) \\
    CPI $\cdot$ surprise model & 0.76 (9) & 0.55 (11) & 0.45 (9) & $-0.37$ (0) \\
    CPI $\cdot$ signal $\times$ vol & 0.77 (12) & 0.74 (12) & 0.61 (11) & $-0.29$ (12) \\
    \midrule
    All $\cdot$ naive & 0.65 (98) & 0.71 (97) & 0.68 (93) & $-2.42$ (82) \\
    All $\cdot$ surprise model & 0.11 (43) & 0.37 (66) & 0.34 (57) & $-0.74$ (2) \\
    All $\cdot$ signal $\times$ vol & 0.56 (101) & 0.37 (100) & 0.29 (96) & $-1.11$ (84) \\
    Combined (CPI model + others $\times$ vol) & 0.40 (98) & 0.23 (101) & 0.16 (95) & $-1.25$ (72) \\
    All $\cdot$ signal $\times$ vol $\cdot$ 2 windows & 0.16 (201) & 0.00 (199) & $-0.12$ (193) & $-2.02$ (168) \\
    \midrule
    All $\cdot$ naive, Sharpe without 2022 & 0.28 & 0.44 & 0.40 & $-2.53$ \\
    \bottomrule
  \end{tabular}
  \begin{tablenotes}\footnotesize
    \item Notes: CPI books use the average of all CPI lines as the signal, so the naive CPI rule
    differs from the headline-only rule of notebook 04. ``All'' is the ten screened families plus CPI.
  \end{tablenotes}
  \end{threeparttable}
\end{table}

\begin{figure}[!htb]
  \centering
  \includegraphics[width=0.85\linewidth]{app_nb05_final_comparison_4markets}
  \caption{CPI only (blue) against CPI plus other surprises (orange: ten screened families; green:
  CPI model plus volatility-scaled others; purple: two windows) and the DTW peer reference (gray):
  trades per year, net Sharpe ratio and net bps per trade, 2019--2026.}
  \label{fig:app_nb05_final}
\end{figure}

\FloatBarrier
\subsection{Surprise and DTW (notebook 06)}

\paragraph{Cumulative P\&L (notebook 06, section 7.3).} Figure~\ref{fig:app_nb06_cum} and
Table~\ref{tab:app_nb06} compare four single-window books on the same releases, walk-forward over
2019--2026 at our base cost: the surprise alone, the DTW direction alone, the surprise kept only when
DTW agrees, and a blend of the two signals. On the broad set of releases the surprise alone is the
best book in \ES\ (0.65) and \NQ\ (0.71); the DTW-only book loses in \ES\ and is flat in \NQ, and
requiring agreement removes about half the trades without raising the Sharpe ratio. \YM\ is the
exception: there DTW alone earns 0.47 and the blend 1.06, with a third of the drawdown, but the
bootstrap probability that the blend is not better than the surprise alone is 6\%, and the same
combination does not help in \ES\ or \NQ\ (Section~\ref{sec:results}). On CPI alone the four books are
close, and most of their cumulative P\&L comes from 2022. In \ZN\ every book loses; the
surprise book loses most because it trades every release and pays a full tick each time.

\begin{table}[!htb]
  \centering\small
  \caption{Surprise, DTW and combined single-window books: net Sharpe ratio, 2019--2026, 2 ticks +
  \$4.50.}\label{tab:app_nb06}
  \begin{tabular}{lcccc}
    \toprule
    Book & \ES & \NQ & \YM & \ZN \\
    \midrule
    \multicolumn{5}{l}{\emph{All screened releases}} \\
    Surprise only & 0.65 & 0.71 & 0.68 & $-2.42$ \\
    DTW only & $-0.39$ & 0.08 & 0.47 & $-1.25$ \\
    Surprise + DTW agree & 0.31 & 0.62 & 0.94 & $-1.57$ \\
    Surprise + DTW blend & 0.41 & 0.67 & 1.06 & $-1.51$ \\
    \midrule
    \multicolumn{5}{l}{\emph{CPI only}} \\
    Surprise only & 0.76 & 0.55 & 0.45 & $-0.37$ \\
    DTW only & 0.32 & 0.81 & 0.41 & $-0.55$ \\
    Surprise + DTW agree & 0.93 & 0.54 & 0.32 & -- \\
    Surprise + DTW blend & 0.69 & 0.72 & 0.51 & $-0.59$ \\
    \bottomrule
  \end{tabular}
  \par\smallskip\footnotesize ``--'': no trades. The \ZN\ CPI surprise book trades once.
\end{table}

\begin{figure}[!htb]
  \centering
  \includegraphics[width=0.82\linewidth]{app_nb06_cumulative_4markets}
  \caption{Cumulative net P\&L of the surprise-only, DTW-only, surprise + DTW agree and surprise + DTW
  blend books on all screened releases (left) and on CPI (right), 2019--2026.}
  \label{fig:app_nb06_cum}
\end{figure}

\paragraph{The ladder: can we reach about 185 trades a year? (notebook 06, section 10).} The
peer DTW strategy trades about 185 times a year on \ES\ because it takes five consecutive
30-minute trades after each release; our books take one trade per release, which caps them at about
100 trades a year. To answer the supervisor's question directly---if we trade five windows after
each release, how many trades do we get, and what happens to the Sharpe ratio?---notebook 06 copies
the ladder structure onto the same releases as the broad surprise book (Table~\ref{tab:app_ladder_design}).
The design was fixed before running. In window $k$ the DTW path is the 31 minutes ending at that
window's entry, matched against earlier releases at the same clock time, and the neighbours'
window-$k$ returns are the forecast; neighbours must have finished window $k$ before the current
release, and the gate is computed separately for each window. The surprise rule holds the sign of
the release's surprise in each window. Every window pays the same round-trip cost. Three books are
compared: a DTW ladder (DTW in all five windows, closest to the peer), a surprise ladder (the surprise
in all five), and the surprise in the first window followed by DTW in windows 2--5. Because the
surprise's effect was already gone in the second window and DTW showed no direction in the first,
we expected the trade count to reach the peer's level while the Sharpe ratio fell toward zero.

\begin{table}[!htb]
  \centering\small
  \caption{Ladder design: five back-to-back windows per release.}\label{tab:app_ladder_design}
  \begin{tabular}{ccc}
    \toprule
    Window & Entry $\to$ exit (bars after the release) & Clock time for an 08:30 release \\
    \midrule
    1 & close of bar $T$ $\to$ $T+29$ & 08:31 $\to$ 09:00 \\
    2 & $T+29 \to T+59$ & 09:00 $\to$ 09:30 \\
    3 & $T+59 \to T+89$ & 09:30 $\to$ 10:00 \\
    4 & $T+89 \to T+119$ & 10:00 $\to$ 10:30 \\
    5 & $T+119 \to T+149$ & 10:30 $\to$ 11:00 \\
    \bottomrule
  \end{tabular}
  \par\smallskip\footnotesize Later windows of an 08:30 release can overlap the first window of a
  10:00 release on the same day; each window is a separate one-contract trade, as in the peer strategy.
\end{table}

The expectation held (Table~\ref{tab:app_ladder_books}, Figure~\ref{fig:app_ladder}). The ladders
reach 300--490 trades a year, well above the peer's 185, but every one earns less than the single
first-window surprise book, and in \ES\ and \ZN\ all of them lose heavily (\ES: $-1.23$, $-1.00$ and
$-0.75$). In \NQ\ the surprise in the first window followed by DTW keeps 0.27, and in \YM\ it is flat.
The extra trades also lower the hit rate below 50\% in \ES\ and \ZN, turn the net return per
trade negative, and raise the maximum drawdown from about $-5\%$ to between $-10\%$ and $-51\%$ in the
equity indices; none of the ladder directions beats random signs at 5\%. More trades per release add
cost, not information. The peer's 185 trades a year with a Sharpe ratio of 0.61 come from its own
release set and setting; on our releases, copying its structure does not reproduce its result.

\begin{table}[!htb]
  \centering\footnotesize
  \caption{Ladder books: trade count against Sharpe ratio (notebook 06, section 10.4), 2019--2026,
  2 ticks + \$4.50.}\label{tab:app_ladder_books}
  \setlength{\tabcolsep}{3pt}
  \begin{tabular}{llcccccccc}
    \toprule
    & Book & Trades/yr & Hit & bps/trade & Total & Sharpe & Ex-2022 & Max DD & $p$ (random) \\
    \midrule
    \ES & Surprise, window 1 only & 98 & 51\% & 1.74 & 12.8\% & 0.65 & 0.28 & $-4.8$\% & 0.00 \\
     & DTW ladder & 314 & 46\% & $-2.06$ & $-48.4$\% & $-1.23$ & $-1.09$ & $-51.2$\% & 0.80 \\
     & Surprise ladder & 490 & 46\% & $-1.15$ & $-42.4$\% & $-1.00$ & $-1.08$ & $-45.6$\% & 0.34 \\
     & Surprise W1 + DTW W2--5 & 348 & 47\% & $-1.10$ & $-28.6$\% & $-0.75$ & $-0.80$ & $-30.1$\% & 0.27 \\
    \midrule
    \NQ & Surprise, window 1 only & 97 & 53\% & 2.41 & 17.5\% & 0.71 & 0.44 & $-5.7$\% & 0.00 \\
     & DTW ladder & 309 & 50\% & $-0.13$ & $-2.9$\% & $-0.07$ & $-0.17$ & $-24.2$\% & 0.28 \\
     & Surprise ladder & 484 & 50\% & $-0.39$ & $-14.3$\% & $-0.29$ & $-0.21$ & $-23.6$\% & 0.43 \\
     & Surprise W1 + DTW W2--5 & 342 & 51\% & 0.50 & 12.9\% & 0.27 & $-0.04$ & $-15.6$\% & 0.06 \\
    \midrule
    \YM & Surprise, window 1 only & 93 & 53\% & 1.73 & 12.0\% & 0.68 & 0.40 & $-4.5$\% & 0.01 \\
     & DTW ladder & 298 & 50\% & $-0.28$ & $-6.1$\% & $-0.23$ & $-0.14$ & $-12.4$\% & 0.18 \\
     & Surprise ladder & 463 & 50\% & $-0.41$ & $-14.2$\% & $-0.34$ & $-0.44$ & $-17.3$\% & 0.15 \\
     & Surprise W1 + DTW W2--5 & 329 & 50\% & 0.00 & 0.1\% & 0.00 & $-0.14$ & $-10.2$\% & 0.06 \\
    \midrule
    \ZN & Surprise, window 1 only & 82 & 36\% & $-3.35$ & $-20.5$\% & $-2.42$ & $-2.53$ & $-20.6$\% & 0.81 \\
     & DTW ladder & 267 & 32\% & $-2.70$ & $-54.1$\% & $-3.92$ & $-4.59$ & $-54.2$\% & 0.10 \\
     & Surprise ladder & 410 & 30\% & $-3.13$ & $-96.1$\% & $-7.65$ & $-8.32$ & $-96.1$\% & 1.00 \\
     & Surprise W1 + DTW W2--5 & 294 & 32\% & $-2.98$ & $-65.7$\% & $-5.35$ & $-6.27$ & $-65.6$\% & 0.89 \\
    \midrule
    & Peer: DTW ladder (slide, \ES, 2018--2026) & 185 & -- & 1.18 & 18.6\% & 0.61 & -- & $-4.9$\% & -- \\
    \bottomrule
  \end{tabular}
  \par\smallskip\footnotesize Net bps per trade; total net return in \% of notional; $p$: share of
  random-sign books with a Sharpe ratio at least as high.
\end{table}

\begin{figure}[!htb]
  \centering
  \includegraphics[width=\linewidth]{app_ladder_cumulative_4markets}
  \caption{Cumulative net P\&L of the five-window ladders (surprise, DTW, and surprise in the first
  window with DTW afterwards) against the single first-window surprise book, 2019--2026.}
  \label{fig:app_ladder}
\end{figure}

\paragraph{Results by window (notebook 06, section 10.3).} Table~\ref{tab:app_ladder_window} scores
each window on its own. The surprise has information only in the first window: its rank correlation
with the next 30-minute return is 0.14 in \ES, 0.09 in \NQ\ and \YM, and then falls to between $-0.06$
and $+0.05$ in windows 2--5, where every surprise window loses. DTW has no consistent information in
any window in \ES\ and \NQ\ (rank correlations of $-0.05$ to $+0.03$); in \YM\ its first window has a
rank correlation of 0.07 and a Sharpe ratio of 0.47, and the later windows fade. In \ZN\ every
window loses under both rules. The per-window Sharpe ratios are shown in Figure~\ref{fig:ladder} of
the main text.

\begin{table}[!htb]
  \centering\small
  \caption{Each ladder window on its own: rank correlation with the next 30-minute return and net
  Sharpe ratio, 2019--2026.}\label{tab:app_ladder_window}
  \setlength{\tabcolsep}{3pt}
  \begin{tabular}{llccccc}
    \toprule
    Market & Rule & Window 1 & Window 2 & Window 3 & Window 4 & Window 5 \\
    \midrule
    \ES & Surprise & 0.14 / 0.65 & 0.01 / $-0.79$ & $-0.06$ / $-0.57$ & 0.00 / $-0.57$ & 0.01 / $-1.16$ \\
        & DTW & 0.01 / $-0.39$ & $-0.02$ / $-0.96$ & 0.01 / $-0.38$ & 0.03 / $-0.22$ & $-0.05$ / $-1.29$ \\
    \NQ & Surprise & 0.09 / 0.71 & $-0.01$ / $-0.23$ & $-0.06$ / $-0.50$ & 0.02 / 0.23 & $-0.04$ / $-0.96$ \\
        & DTW & 0.03 / 0.08 & $-0.01$ / 0.01 & $-0.01$ / $-0.23$ & 0.02 / 0.15 & $-0.01$ / $-0.18$ \\
    \YM & Surprise & 0.09 / 0.68 & 0.05 / $-0.15$ & 0.02 / $-0.18$ & $-0.06$ / $-0.80$ & 0.04 / $-0.30$ \\
        & DTW & 0.07 / 0.47 & 0.03 / $-0.19$ & 0.05 / 0.15 & 0.01 / 0.03 & $-0.06$ / $-1.03$ \\
    \ZN & Surprise & $-0.03$ / $-2.42$ & $-0.01$ / $-3.08$ & 0.01 / $-3.15$ & $-0.04$ / $-3.90$ & 0.04 / $-3.56$ \\
        & DTW & 0.04 / $-1.25$ & $-0.02$ / $-1.94$ & $-0.02$ / $-2.34$ & 0.01 / $-1.62$ & $-0.01$ / $-2.92$ \\
    \bottomrule
  \end{tabular}
  \par\smallskip\footnotesize Each cell: rank correlation / net Sharpe ratio.
\end{table}

\FloatBarrier
\subsection{Multi-market surprise book (notebook 07)}

\paragraph{Performance by market (notebook 07, section 5).} Figure~\ref{fig:app_perf_market} and
Table~\ref{tab:app_perf_market} run the same five books, with the same code and constants, on each
market. The pattern is the same in the three equity indices and opposite in \ZN. The broad naive
surprise book earns 0.65 in \ES, 0.71 in \NQ\ and 0.68 in \YM, with random-sign $p$-values of 0.006 or
below; the CPI model 0.45--0.76; and keeping CPI trades only when DTW agrees 0.32--0.93, on six or
seven trades a year. The DTW-only book loses in \ES, is flat in \NQ\ and earns 0.47 in \YM. Gross
against net Sharpe ratios show what costs take: about 0.5 in \ES, 0.2 in \NQ\ and 0.3 in \YM\ for the
broad book, but 2.1 in \ZN, where the naive book goes from $-0.30$ before costs to $-2.42$ after them. Over the
cost grid the broad book's Sharpe ratio falls from 1.09 to 0.20 in \ES\ as slippage rises from zero
to four ticks, but only from 0.82 to 0.60 in \NQ\ and from 0.91 to 0.45 in \YM, because the \NQ\ and
\YM\ ticks are smaller relative to their moves.

\begin{table}[!htb]
  \centering\small
  \caption{Performance by market, 2019--2026, 2 ticks + \$4.50: net Sharpe ratio (gross in
  parentheses) and trades per year.}\label{tab:app_perf_market}
  \begin{threeparttable}
  \setlength{\tabcolsep}{4pt}
  \begin{tabular}{lcccc}
    \toprule
    Book & \ES & \NQ & \YM & \ZN \\
    \midrule
    All $\cdot$ naive & 0.65 (1.17) & 0.71 (0.87) & 0.68 (1.01) & $-2.42$ ($-0.30$) \\
    \quad trades / yr; $p$ (random sign) & 98; 0.000 & 97; 0.005 & 93; 0.006 & 82; 0.81 \\
    All $\cdot$ DTW only & $-0.39$ ($-0.01$) & 0.08 (0.21) & 0.47 (0.79) & $-1.25$ (0.47) \\
    All $\cdot$ surprise + DTW agree & 0.31 (0.70) & 0.62 (0.73) & 0.94 (1.21) & $-1.57$ ($-0.12$) \\
    CPI $\cdot$ surprise model & 0.76 (0.88) & 0.55 (0.60) & 0.45 (0.53) & $-0.37$ (1 trade) \\
    CPI $\cdot$ surprise + DTW agree & 0.93 (1.02) & 0.54 (0.57) & 0.32 (0.38) & -- \\
    \midrule
    All $\cdot$ naive, Sharpe without 2022 & 0.28 & 0.44 & 0.40 & $-2.53$ \\
    \midrule
    \multicolumn{5}{l}{\emph{All $\cdot$ naive, net Sharpe by round-trip slippage (plus \$4.50)}} \\
    0 ticks & 1.09 & 0.82 & 0.91 & $-0.56$ \\
    1 tick  & 0.87 & 0.77 & 0.80 & $-1.49$ \\
    2 ticks & 0.65 & 0.71 & 0.68 & $-2.42$ \\
    4 ticks & 0.20 & 0.60 & 0.45 & $-4.24$ \\
    \bottomrule
  \end{tabular}
  \begin{tablenotes}\footnotesize
    \item Notes: Monthly Sharpe ratios; ``--'': no trades. The \ZN\ CPI model trades about once over
    the sample, because its predicted drift rarely exceeds the cost.
  \end{tablenotes}
  \end{threeparttable}
\end{table}

\begin{figure}[!htb]
  \centering
  \includegraphics[width=0.85\linewidth]{app_nb07b_sharpe_by_market}
  \caption{Net Sharpe ratio by market and strategy, 2019--2026, 2 ticks + \$4.50 (notebook 07,
  section 5).}
  \label{fig:app_perf_market}
\end{figure}

\paragraph{Portfolios: all markets, equities and \ES\ + \NQ\ (notebook 07, sections 5--6).}
Table~\ref{tab:app_portfolios07} and Figure~\ref{fig:app_multi} combine the per-market books with
equal-risk weights fixed on 2016--2018 (\ES\ 1.00, \NQ\ 0.71, \YM\ 1.02, \ZN\ 2.67). Adding \YM\ to \ES\ and
\NQ\ improves the multi-release equity book: the Sharpe ratio rises from 0.72 to 0.78, total net P\&L
from 25\% to 38\%, trades from 195 to 288 a year, and the 90\% confidence interval moves fully above
zero (0.10 to 1.50). The Sharpe ratio without 2022 also rises (0.38 to 0.43). The cost is a larger
drawdown ($-9.2\%$ against $-7.4\%$), because the book is bigger. \YM\ helps because its monthly
returns correlate only 0.66 with \ES's and 0.43 with \NQ's. For the CPI-only model \YM\ adds little
(0.69 to 0.63), because the CPI books of the three indices are highly correlated. \ZN\ breaks every
portfolio it enters: equal-risk weighting gives it the largest weight (2.67, because its moves are
small), its multi-release book loses ($-2.42$), and the four-market naive book is negative ($-0.34$,
drawdown $-23\%$). \ZN\ should therefore be excluded from the traded book. The four-market CPI
portfolio still earns 0.61 because the CPI model almost never trades \ZN.

\begin{table}[!htb]
  \centering\small
  \caption{Equal-risk surprise portfolios across markets, 2019--2026, 2 ticks + \$4.50 (weights fixed
  on 2016--2018: \ES\ 1.00, \NQ\ 0.71, \YM\ 1.02, \ZN\ 2.67).}\label{tab:app_portfolios07}
  \setlength{\tabcolsep}{3pt}
  \begin{tabular}{lccccccc}
    \toprule
    Portfolio & Trades/yr & Sharpe & 90\% CI & Ex-2022 & Total & Max DD & $p$ (random) \\
    \midrule
    \ES+\NQ, all $\cdot$ naive & 195 & 0.72 & [$-0.01$, 1.43] & 0.38 & 25.2\% & $-7.4\%$ & 0.000 \\
    \textbf{Equities (\ES+\NQ+\YM), all $\cdot$ naive} & \textbf{288} & \textbf{0.78} & \textbf{[0.10, 1.50]} & \textbf{0.43} & \textbf{37.5\%} & $-9.2\%$ & \textbf{0.000} \\
    \ES+\NQ, CPI model & 20 & 0.69 & [$-0.01$, 1.39] & 0.02 & 9.9\% & $-3.3\%$ & 0.023 \\
    Equities (\ES+\NQ+\YM), CPI model & 29 & 0.63 & [$-0.10$, 1.37] & $-0.04$ & 13.0\% & $-5.3\%$ & 0.020 \\
    All four markets, all $\cdot$ naive & 370 & $-0.34$ & [$-1.08$, 0.40] & $-0.74$ & $-17.4\%$ & $-23.4\%$ & 0.000 \\
    All four markets, CPI model & 29 & 0.61 & [$-0.13$, 1.35] & $-0.07$ & 12.6\% & $-5.8\%$ & 0.018 \\
    \midrule
    All four markets, DTW only & 244 & $-0.64$ & [$-1.28$, $-0.01$] & $-0.62$ & $-23.7\%$ & $-29.1\%$ & 0.028 \\
    All four markets, surprise + DTW agree & 195 & $-0.07$ & [$-0.74$, 0.61] & $-0.29$ & $-2.6\%$ & $-15.6\%$ & 0.001 \\
    All four markets, CPI + DTW agree & 19 & 0.67 & [$-0.01$, 1.31] & $-0.01$ & 11.9\% & $-3.4\%$ & 0.010 \\
    \bottomrule
  \end{tabular}
  \par\smallskip\footnotesize Total net return in \% of \ES-equivalent notional. For the four-market
  naive book, the random-sign $p$ of 0.000 says its direction is better than random; the loss comes
  from \ZN's costs, which random-sign books pay too.
\end{table}

\begin{figure}[!htb]
  \centering
  \includegraphics[width=0.8\linewidth]{app_nb07b_portfolio_cumulative}
  \caption{Cumulative net P\&L of the risk-weighted surprise books (broad naive book) by market and
  of the two-, three- and four-market portfolios, 2019--2026.}
  \label{fig:app_multi}
\end{figure}

\paragraph{Final comparison (notebook 07, section 7).} Table~\ref{tab:app_final07} brings every
market and portfolio book of notebook 07 together. It summarizes the notebook's conclusion: the
consensus surprise is the source of direction in all three equity indices, the broad naive book is
the most robust single book (positive without 2022 and significant against random signs in each
index), the three-index portfolio is the best overall, and \ZN\ should not be traded at these costs.

\begin{table}[!htb]
  \centering\footnotesize
  \caption{Final comparison of the multi-market surprise books, 2019--2026, 2 ticks + \$4.50
  (notebook 07, section 7).}\label{tab:app_final07}
  \setlength{\tabcolsep}{4pt}
  \begin{tabular}{lcccccc}
    \toprule
    Book & Trades/yr & bps/trade & Sharpe & Ex-2022 & Max DD & $p$ (random) \\
    \midrule
    \ES\ $\cdot$ CPI $\cdot$ surprise model & 9 & 8.23 & 0.76 & 0.05 & $-2.0\%$ & 0.006 \\
    \ES\ $\cdot$ All $\cdot$ naive & 98 & 1.74 & 0.65 & 0.28 & $-4.8\%$ & 0.000 \\
    \ES\ $\cdot$ All $\cdot$ surprise + DTW agree & 53 & 1.14 & 0.31 & 0.07 & $-4.6\%$ & 0.030 \\
    \ES\ $\cdot$ All $\cdot$ DTW only & 65 & $-1.45$ & $-0.39$ & $-0.32$ & $-11.0\%$ & 0.479 \\
    \midrule
    \NQ\ $\cdot$ CPI $\cdot$ surprise model & 11 & 7.12 & 0.55 & $-0.01$ & $-2.7\%$ & 0.047 \\
    \NQ\ $\cdot$ All $\cdot$ naive & 97 & 2.41 & 0.71 & 0.44 & $-5.7\%$ & 0.005 \\
    \NQ\ $\cdot$ All $\cdot$ surprise + DTW agree & 51 & 3.11 & 0.62 & 0.59 & $-6.7\%$ & 0.021 \\
    \NQ\ $\cdot$ All $\cdot$ DTW only & 63 & 0.34 & 0.08 & 0.17 & $-8.3\%$ & 0.271 \\
    \midrule
    \YM\ $\cdot$ CPI $\cdot$ surprise model & 9 & 4.67 & 0.45 & $-0.15$ & $-2.0\%$ & 0.060 \\
    \YM\ $\cdot$ All $\cdot$ naive & 93 & 1.73 & 0.68 & 0.40 & $-4.5\%$ & 0.006 \\
    \YM\ $\cdot$ All $\cdot$ surprise + DTW agree & 49 & 2.94 & 0.94 & 0.94 & $-1.7\%$ & 0.002 \\
    \YM\ $\cdot$ All $\cdot$ DTW only & 61 & 1.27 & 0.47 & 0.58 & $-4.0\%$ & 0.024 \\
    \midrule
    \ZN\ $\cdot$ CPI $\cdot$ surprise model & 0 & $-17.03$ & $-0.37$ & $-0.39$ & $-0.2\%$ & 1.000 \\
    \ZN\ $\cdot$ All $\cdot$ naive & 82 & $-3.35$ & $-2.42$ & $-2.53$ & $-20.6\%$ & 0.812 \\
    \ZN\ $\cdot$ All $\cdot$ surprise + DTW agree & 42 & $-3.18$ & $-1.57$ & $-1.89$ & $-10.2\%$ & 0.645 \\
    \ZN\ $\cdot$ All $\cdot$ DTW only & 55 & $-2.15$ & $-1.25$ & $-1.71$ & $-9.1\%$ & 0.164 \\
    \midrule
    2-market (\ES+\NQ) $\cdot$ All $\cdot$ naive & 195 & 1.73 & 0.72 & 0.38 & $-7.4\%$ & 0.000 \\
    2-market (\ES+\NQ) $\cdot$ CPI $\cdot$ surprise model & 20 & 6.45 & 0.69 & 0.02 & $-3.3\%$ & 0.023 \\
    \midrule
    3-equity (\ES+\NQ+\YM) $\cdot$ All $\cdot$ naive & 288 & 1.74 & 0.78 & 0.43 & $-9.2\%$ & 0.000 \\
    3-equity (\ES+\NQ+\YM) $\cdot$ CPI $\cdot$ surprise model & 29 & 5.94 & 0.63 & $-0.04$ & $-5.3\%$ & 0.020 \\
    \midrule
    4-market $\cdot$ CPI $\cdot$ surprise model & 29 & 5.71 & 0.61 & $-0.07$ & $-5.8\%$ & 0.018 \\
    4-market $\cdot$ All $\cdot$ naive & 370 & $-0.63$ & $-0.34$ & $-0.74$ & $-23.4\%$ & 0.000 \\
    4-market $\cdot$ All $\cdot$ DTW only & 244 & $-1.30$ & $-0.64$ & $-0.62$ & $-29.1\%$ & 0.028 \\
    4-market $\cdot$ All $\cdot$ surprise + DTW agree & 195 & $-0.18$ & $-0.07$ & $-0.29$ & $-15.6\%$ & 0.001 \\
    4-market $\cdot$ CPI $\cdot$ surprise + DTW agree & 19 & 8.17 & 0.67 & $-0.01$ & $-3.4\%$ & 0.010 \\
    \bottomrule
  \end{tabular}
  \par\smallskip\footnotesize Net bps per trade; monthly Sharpe ratios; $p$: share of random-sign
  books with a Sharpe ratio at least as high.
\end{table}



